% Topic presentation. Template format: 11pt, 1-inch margins, single column.
\documentclass[11pt]{article}
\usepackage[utf8]{inputenc}
\usepackage[T1]{fontenc}
\usepackage[spanish,es-nodecimaldot]{babel}
\usepackage[margin=1in]{geometry}
\usepackage{amsmath,amssymb,booktabs}
\usepackage[round]{natbib}
\usepackage{xcolor}
\usepackage[colorlinks=true,allcolors=blue!50!black]{hyperref}
\newcommand{\E}{\mathbb E}
\newcommand{\Var}{\operatorname{Var}}
\title{\textbf{¿Cuándo la transparencia mejora la incorporación\\
de fundamentos al precio?}}
\author{David Julio Olano Silva Nevado\\
\small Universidad del Pacífico\\
\small\url{https://github.com/davidolano03/ai-project}}
\date{Presentación del tema: 9 de octubre de 2026}
\begin{document}
\maketitle

\section{La pregunta y la ruta}
\textbf{¿Bajo qué condiciones una mayor transparencia informativa aumenta la
sensibilidad esperada del precio a los fundamentos cuando algunos
inversionistas están afectados por sentimiento?} Seguiré la \textbf{ruta A}:
corregiré la representación de información pública en \citet{zhang2023},
conservaré sus cuatro grupos y volveré a derivar el resultado afectado.
El objetivo es distinguir noticias sobre el pago futuro de expectativas
sesgadas. Una respuesta más informativa del precio significa que refleja
mejor tanto las buenas como las malas noticias; no que necesariamente suba.

\section{El punto de partida y el supuesto que cambiaré}
La referencia exacta es la \textbf{proposición 5(ii), p.~16}, que afirma que
todo grado de transparencia $0\leq\mu\leq1$ aumenta la sensibilidad del precio
a la información respecto del modelo básico. Cambiaré las ecuaciones
(38)--(41), pp.~15--16. Allí los particulares observan
$\kappa=\mu\theta$, donde $\theta$ es la noticia fundamental y $\kappa$ la
información pública, y su varianza subjetiva incluye
\begin{equation}\label{eq:printed}
  v_3^I=\sigma_\epsilon^2+\mu^2\sigma_\theta^2.
\end{equation}
Esta expresión aumenta con la parte divulgada y, en $\mu=0$, elimina la
incertidumbre sobre una noticia que el particular todavía desconoce.
Además, observar exactamente $\kappa=\mu\theta$, con $\mu>0$ conocido,
permite recuperar $\theta=\kappa/\mu$. La información parcial necesita una
especificación distinta de una señal escalar multiplicada por una constante.

\textbf{Evidencia preliminar.} Con participaciones $\lambda_i=1/4$,
$R=\gamma=1$, sentimientos y medias cero, y todas las varianzas del modelo
básico iguales a uno, las fórmulas (23) y (44) publicadas dan
\[
 K_0=\frac9{14}\simeq0{,}642857,
 \qquad K_I(0{,}1)=\frac{66943}{121303}\simeq0{,}551866<K_0.
\]
Lean verifica este contraejemplo con transparencia positiva. Refuta la
afirmación universal \emph{con las fórmulas impresas}; no prueba que la
transparencia real perjudique los mercados. La corrección propuesta se
concentra en esta inconsistencia, no en invalidar todo el artículo.

\newpage
\subsection*{Información parcial coherente y problema del agente}
Descompondré la noticia en componentes gaussianos independientes, de media
cero: $\theta=\kappa+\xi$. Todos observan $\kappa$ y las instituciones también
observan $\xi$. Si $s=\sigma_\theta^2>0$ es la varianza total, impondré
\begin{equation}\label{eq:information}
 \Var(\kappa)=\mu s,\qquad \Var(\xi)=(1-\mu)s,\qquad
 \E[\theta\mid\kappa]=\kappa,\quad
 \Var(\theta\mid\kappa)=(1-\mu)s.
\end{equation}
Aquí $\mu$ mide la fracción de \emph{varianza informativa} que se divulga,
no una fracción observable del valor realizado de $\theta$. Más transparencia
reduce la incertidumbre que queda. En $\mu=0$ recupero el modelo básico;
en $\mu=1$ todos conocen $\theta$, pero persiste la sorpresa $\epsilon$.

El pago final es $F=P_0+\theta+\epsilon$, con $P_0=P_{t-1}$ conocido y
$\epsilon$ de media cero e independiente de los componentes informativos.
$P$ es el precio actual y $R=1+r_f>0$ el factor bruto seguro. Mantendré
los grupos $i=1,2,3,4$ y sus participaciones $\lambda_i>0$,
$\sum_i\lambda_i=1$. Denoto $m_i$ y $v_i$ como media y varianza subjetivas
del pago; $\delta_j=\rho_j^*$ como sesgo esperado y
$\tau_j^2=\sigma_{\rho_j}^2$ como incertidumbre subjetiva adicional.
Distingo estos parámetros del estado emocional observado $\rho_j$.
La aversión del grupo $i$ es $g_i$: $g_1=g_3=\gamma$,
$g_2=\gamma e^{-b\rho_1}$ y $g_4=\gamma e^{-b\rho_2}$, con $b,\gamma>0$.

\begin{center}\small
\begin{tabular}{@{}llll@{}}
\toprule
Grupo & $m_i$ & $v_i$ & Aversión\\
\midrule
1. Institución racional & $P_0+\theta$ & $\sigma_\epsilon^2$ & $\gamma$\\
2. Institución con sentimiento & $P_0+\theta+\delta_1$ & $\sigma_\epsilon^2+\tau_1^2$ & $g_2$\\
3. Particular racional & $P_0+\kappa$ & $\sigma_\epsilon^2+(1-\mu)s$ & $\gamma$\\
4. Particular con sentimiento & $P_0+\kappa+\delta_2$ & $\sigma_\epsilon^2+(1-\mu)s+\tau_2^2$ & $g_4$\\
\bottomrule
\end{tabular}
\end{center}
Las instituciones conocen toda la noticia; los particulares sólo su parte
pública. Los sesgos y las varianzas adicionales son subjetivos: no añado
sentimiento al pago objetivo de la empresa.

Con riqueza final $W_i=C_{0,i}+\omega_i(F-RP)$, cada agente elige
su posición $\omega_i\in\mathbb R$ para maximizar utilidad CARA esperada.
Bajo sus creencias normales, el problema equivalente y la FOC esperada son
\begin{equation}\label{eq:foc}
 \max_{\omega_i\in\mathbb R}
 \left\{C_{0,i}+\omega_i(m_i-RP)-\frac{g_iv_i}{2}\omega_i^2\right\},
 \qquad m_i-RP-g_iv_i\omega_i=0.
\end{equation}
$C_{0,i}$ es riqueza independiente de la decisión. El término lineal es
ganancia esperada; el cuadrático penaliza riesgo. Como $-g_iv_i<0$,
$\omega_i^*=(m_i-RP)/(g_iv_i)$ es el máximo único. Cambiaré las medias y
varianzas de los particulares, no la forma de la FOC. Para el grupo 3, el
coeficiente de $\omega_3$ pasa a ser
$\gamma[\sigma_\epsilon^2+(1-\mu)s]$. Con precio y expectativas iguales,
menos riesgo amplía la posición: más compra si es positiva y más venta corta
si es negativa. No impondré restricciones de ventas cortas en esta entrega.

\newpage
\section{El resultado que espero y su interpretación}
La oferta del activo es una unidad. Con
$h_i=\lambda_i/(g_iv_i)$ y $H=\sum_i h_i$, el equilibrio requiere
\begin{equation}\label{eq:price}
 \sum_i\lambda_i\omega_i^*=1,
 \qquad P^*=\frac{\sum_i h_im_i-1}{RH}.
\end{equation}
$h_i$ es el peso efectivo de demanda: crece con la participación y disminuye
con aversión e incertidumbre. Una mayor relevancia de un grupo aumenta el
precio \emph{si su valoración $m_i$ supera $RP$}; con una valoración menor,
lo reduce. El término $1/H$ refleja la compensación por absorber la oferta
riesgosa y no debe confundirse con un sesgo de expectativas.

Mi comparación promediará las señales públicas compatibles con un mismo
fundamento:
\[
 \bar P(\theta,\mu)=\E[P^*\mid\theta,\rho_1,\rho_2].
\]
Por normalidad e independencia, $\E[\kappa\mid\theta]=\mu\theta$.
Los estados sentimentales serán exógenos, independientes de $\kappa,\xi$
y fijos al comparar $\mu$. Definiendo
\begin{align*}
 a&=h_1+h_2>0,\qquad n_3=\lambda_3/\gamma>0,\qquad n_4=\lambda_4/g_4>0,\\
 d_3(\mu)&=\sigma_\epsilon^2+(1-\mu)s,\qquad
 d_4(\mu)=\sigma_\epsilon^2+\tau_2^2+(1-\mu)s,\\
 B(\mu)&=\frac{n_3}{d_3(\mu)}+\frac{n_4}{d_4(\mu)},
\end{align*}
obtendré como medida de incorporación de fundamentos
\begin{equation}\label{eq:K}
 \mathcal K(\mu)=\frac{\partial\bar P}{\partial\theta}
       =\frac{a+\mu B(\mu)}{R[a+B(\mu)]}.
\end{equation}
El numerador suma la respuesta directa de instituciones y la respuesta
esperada de particulares. El denominador agrega todos los pesos de demanda
y descuenta en el tiempo. Es una sensibilidad del \emph{precio esperado};
no la derivada del precio realizado manteniendo fija la señal $\kappa$.

\textbf{Resultado objetivo.} Para $0\leq x<y\leq1$,
$\mathcal K(y)>\mathcal K(x)$ y $\mathcal K(1)=1/R$, si las participaciones,
sesgos, estados de sentimiento, aversiones de referencia y varianzas totales
permanecen fijos, con $\sigma_\epsilon^2>0$, $s>0$, $\tau_j^2\geq0$ y
$R>0$. La parte algebraica de este resultado ya cuenta con una prueba
preliminar en Lean para los cuatro grupos. La derivación esperada es
\begin{equation}\label{eq:derivative}
 \mathcal K'(\mu)=
 \frac{B^2+a\left[\dfrac{n_3\sigma_\epsilon^2}{d_3^2}
                 +\dfrac{n_4(\sigma_\epsilon^2+\tau_2^2)}{d_4^2}\right]}
      {R(a+B)^2}>0.
\end{equation}
La transparencia permite que los particulares usen la noticia y reduce su
riesgo residual. Una buena noticia mueve más el precio hacia arriba y una
mala noticia más hacia abajo. No afirmaré que el precio siempre aumenta,
que desaparecen los sesgos ni que mejora necesariamente el bienestar.
La corrección sustituye la comparación incondicional publicada por un
resultado con medida y condiciones explícitas.

\newpage
\section{Por qué esta corrección requiere trabajo propio}
Revisé la sección 6 completa y el apéndice, pp.~17--18, del artículo
publicado: ya incluyen aprendizaje y transparencia, y la derivación de
la proposición 5 conserva la parametrización cuestionada. Mi aporte será
cambiar la estructura informativa, distinguir varianza observada de residual
y volver a probar el resultado para cuatro grupos.

La búsqueda documentada por título exacto y DOI, incluida la búsqueda de
versiones de trabajo y correcciones, no identificó una versión alternativa
accesible de Zhang ni una reparación de esas ecuaciones. Revisé además
\citet{das2024}, un artículo que cita Zhang y estudia empíricamente miedo,
información asimétrica y riesgo de caídas bursátiles. Su versión de trabajo
de Canterbury 1/2024, pp.~23--24 del texto, usa Zhang como motivación de
sentimiento e información; no ofrece esta rederivación del modelo.
El registro \texttt{proposal/literature-search.md} distingue texto revisado,
metadatos y accesos pendientes. Esta revisión delimitada respalda una
corrección específica, no una afirmación de novedad universal de la idea
de información pública. Ampliaré la revisión de artículos citantes antes
del manuscrito final.

\section{Plan, verificación y riesgos}
\textbf{Prueba.} Derivaré las medias y varianzas condicionales, las cuatro
FOC, el equilibrio y el efecto de $\mu$, nombrando cada supuesto.
Separaré el precio realizado de su respuesta esperada. La evidencia
preliminar incluye el contraejemplo y la desigualdad algebraica de
$\mathcal K$; el puente probabilístico todavía no está formalizado en Lean.

\textbf{Simulación.} \texttt{python code/verify.py} reproduce las fracciones,
el gráfico y los extremos, verifica las FOC y el vaciado de mercado, y
comprueba el signo sobre una malla de parámetros admisibles. Estas
simulaciones ilustran y detectan errores; no sustituyen una prueba general.

\textbf{Lean.} Mapearé los resultados a declaraciones con hipótesis
explícitas y sin \texttt{sorry}. La entrega incluye 15 comprobaciones de
la auditoría inicial y 4 de la corrección con grupos particulares distintos.
El procedimiento de AppliedModelingLib del curso se realizará sobre el
commit fijado del manuscrito final. Esta auditoría preliminar tiene un
alcance distinto de la formalización integral prevista.

\textbf{Riesgo principal.} El resultado puede perderse si la divulgación
cambia endógenamente el sentimiento, las participaciones o la aversión.
Mantendré esas variables fijas en el resultado principal y trataré su
variación como un escenario adicional con condiciones por derivar. También
delimitaré la falta de inferencia adicional desde precios, manteniendo
la restricción informativa del punto de partida. Elaboraré personalmente
el apéndice manuscrito; una guía mecanografiada no lo reemplazará.

\bibliographystyle{plainnat}
\bibliography{../paper/references}
\end{document}
